% --------------------------------------------------------------------------------------- %
%                   UOC - CSD - HY119 - Spring 2025 - Ass1 - csd5748 - :)                 %
% --------------------------------------------------------------------------------------- %

\documentclass[fleqn]{article}
\setlength{\mathindent}{0pt}
\usepackage[utf8]{inputenc} 
\usepackage[english, greek]{babel}
\usepackage{subcaption}
\usepackage{hyperref}
\usepackage{xcolor}
\hypersetup
{
    hidelinks,  
    colorlinks=true,
    linkcolor=blue,
    urlcolor=blue 
}
\captionsetup[subfigure]{labelformat=empty}
\usepackage{amsmath}
\usepackage{amsfonts}
\usepackage{amssymb}
\usepackage{array}
\usepackage{graphicx}
\usepackage[a4paper, left=0.7cm, right=0.7cm, top=1.3cm, bottom=2.6cm]{geometry}
\usepackage[absolute, overlay]{textpos}
\usepackage{listings}
\usepackage{xcolor}
\setlength{\jot}{0pt}
\lstdefinestyle{mypython}
{
    language=Python,
    basicstyle=\ttfamily\small,
    keywordstyle=\color{blue},
    stringstyle=\color{red},
    commentstyle=\color{green!50!black},
    numbers=left,
    numberstyle=\tiny\color{gray},
    stepnumber=1,
    frame=single,
    backgroundcolor=\color{gray!10},
    breaklines=true
}
\begin{document}
\selectlanguage{greek}

% ---------------------------------------------------------------------------------------

\begin{center}
\boxed{\textlatin{csd5748}}
\end{center}
\vspace{-0.735cm}
\begin{center}
\setlength{\fboxsep}{20pt}
\noindent{\boxed{{{\textbf{\Large{ΗΥ-119 - Γραμμική Άλγεβρα -  Άσκηση 1.}}}}}}
\end{center}

% ---------------------------------------------------------------------------------------

\vspace{0.999cm}
{{\textbf{·  \underline{1\textsuperscript{η} Άσκηση:}}}}
\vspace{0.54321cm}

\[
\textbf{Έχουμε:} \hspace{0.54321cm}
\left\{
\begin{array}{l}
2u - 3v + 0w = 3 \\
4u + 5v + 1w = 7 \\
2u - 1v - 3w = 5
\end{array}
\right.
\hspace{0.321cm} \overset{\textbf{Π}_{3\times4}}{\implies} \hspace{0.321cm}
\begin{bmatrix}
& \boxed{2} & -3 & 0 & | & 3 & \\
& 4 & 5 & 1 & | & 7 & \\
& 2 & -1 & -3 & | & 5 &
\end{bmatrix}
\hspace{0.321cm} \overset{r_2=r_2-\lambda r_1}{\underset{\lambda=\frac{4}{2}=2}{\implies}} \hspace{0.321cm}
\begin{bmatrix}
& \boxed{2} & -3 & 0  & 3 & \\
& \underline{ 0 } & 11 & 1 &  1 & \\
& 2 & -1 & -3 & 5 &
\end{bmatrix}
\]

\[
\overset{r_3=r_3-\lambda r_1}{\underset{\lambda=\frac{2}{2}=1}{\implies}} \hspace{0.321cm}
\begin{bmatrix}
& 2 & -3 & 0  & 3 & \\
& \underline{ 0 } & \boxed{11} & 1 &  1 & \\
& \underline{0} & 2 & -3 & 2 &
\end{bmatrix}
\hspace{0.321cm} \overset{r_3=r_3-\lambda r_2}{\underset{\lambda=2/11}{\implies}} \hspace{0.321cm}
\begin{bmatrix}
& 2 & -3 & 0  & 3 & \\
& \underline{ 0 } & 11 & 1 &  1 & \\
& \underline{0} & \underline{0} & \boxed{-\frac{35}{11}} & \frac{20}{11} &
\end{bmatrix}
\hspace{0.321cm} {\implies} \hspace{0.321cm}
\begin{bmatrix}
& 2 & -3 & 0 & | & 3 & \\
& 0 & 11 & 1 & | & 1 & \\
& 0 & 0 & {-35/11} & | & 20/11 &
\end{bmatrix}
\hspace{0.234cm}
\]

\[
\hspace{0.321cm} {\implies} \hspace{0.321cm}
\left\{
\begin{array}{l}
2u - 3v + 0w = 3 \\
0u + 11v + 1w = 1 \\
0u - 0v - \frac{35}{11}w = \frac{20}{11}
\end{array}
\right.
\hspace{0.321cm} \overset{r_3:}{\implies} \hspace{0.321cm}
\left\{
\begin{array}{l}
2u - 3v + 0w = 3 \\
0u + 11v + 1w = 1 \\
- \frac{35}{11}w = \frac{20}{11}
\end{array}
\right.
\hspace{0.321cm} {\implies} \hspace{0.321cm}
\left\{
\begin{array}{l}
2u - 3v + 0w = 3 \\
0u + 11v + 1w = 1 \\
w = -\frac{4}{7}
\end{array}
\right.
\hspace{0.321cm} {\implies} \hspace{0.321cm}
\]

\[
\hspace{0.321cm} \overset{r_2:}{\implies} \hspace{0.321cm}
\left\{
\begin{array}{l}
2u - 3v + 0w = 3 \\
11v + 1 \cdot (-\frac{4}{7}) = 1 \\
w = -\frac{4}{7}
\end{array}
\right.
\hspace{0.321cm} {\implies} \hspace{0.321cm}
\left\{
\begin{array}{l}
2u - 3v + 0w = 3 \\
v = \frac{1}{7} \\
w = -\frac{4}{7}
\end{array}
\right.
\hspace{0.321cm} \overset{r_1:}{\implies} \hspace{0.321cm}
\left\{
\begin{array}{l}
2u - 3 \cdot (\frac{1}{7})= 3 \\
v = \frac{1}{7} \\
w = -\frac{4}{7}
\end{array}
\right.
\hspace{0.321cm} \overset{}{\implies} \hspace{0.321cm}
\]

\[
\hspace{0.321cm} \overset{... \hspace{0.150cm}}{\implies} \hspace{0.321cm}
\left\{
\hspace{0.123456789cm} \boxed
{
\begin{array}{l}
u = 12 / 7 \\
v = 1 / 7 \\
w = - 4 / 7
\end{array}
}
\right.
\]


% ---------------------------------------------------------------------------------------


\vspace{2.000cm}
{{\textbf{·  \underline{2\textsuperscript{η} Άσκηση:}}}}
\vspace{0.321cm}

\[
i) \hspace{0.321cm} \textbf{Έχουμε:} \hspace{0.54321cm}
\begin{bmatrix}
& 4 & 0 & 1 & \\
& 0 & 1 & 0 & \\
& 4 & 0 & 0 &
\end{bmatrix}_{3 \times 3}^A
\cdot \hspace{0.35cm}
\begin{bmatrix}
& 3 & \\
& 4 & \\
& 5 &
\end{bmatrix}_{3 \times 1}^B
\hspace{0.321cm} \overset{*}{=} \hspace{0.321cm}
\begin{bmatrix}
& 
\begin{bmatrix}
& 4 & 0 & 1 &
\end{bmatrix}
\cdot
\begin{bmatrix}
& 3 & \\
& 4 & \\
& 5 &
\end{bmatrix} 
& \\
& 
\begin{bmatrix}
& 0 & 1 & 0 &
\end{bmatrix}
\cdot
\begin{bmatrix}
& 3 & \\
& 4 & \\
& 5 &
\end{bmatrix} 
& \\
& 
\begin{bmatrix}
& 4 & 0 & 0 &
\end{bmatrix}
\cdot
\begin{bmatrix}
& 3 & \\
& 4 & \\
& 5 &
\end{bmatrix} 
&
\end{bmatrix}
\hspace{0.321cm} \overset{}{=}
\]

\[
\hspace{0.321cm} \overset{}{=} \hspace{0.321cm}
\begin{bmatrix}
& (4)(3) + (0)(4) + (1)(5) & \\
& (0)(3) + (1)(4) + (0)(5) & \\
& (4)(3) + (0)(4) + (0)(5) &
\end{bmatrix}
\hspace{0.321cm} {=} \hspace{0.321cm}
\boxed
{
\begin{bmatrix}
& 21 & \\
& 4 & \\
& 12 &
\end{bmatrix}_{3 \times 1}^C
}
\]


\newpage


\[
ii) \hspace{0.321cm} \textbf{Έχουμε:} \hspace{0.54321cm}
\begin{bmatrix}
& 1 & 0 & 0 & \\
& 0 & 1 & 0 & \\
& 0 & 0 & 1 &
\end{bmatrix}_{3 \times 3}^I
\cdot \hspace{0.35cm}
\begin{bmatrix}
& 5 & \\
& -2 & \\
& 3 &
\end{bmatrix}_{3 \times 1}^D
\hspace{0.321cm} \overset{*}{=} \hspace{0.321cm}
\begin{bmatrix}
& 
\begin{bmatrix}
& 1 & 0 & 0 &
\end{bmatrix}
\cdot
\begin{bmatrix}
& 5 & \\
& -2 & \\
& 3 &
\end{bmatrix} 
& \\
& 
\begin{bmatrix}
& 0 & 1 & 0 &
\end{bmatrix}
\cdot
\begin{bmatrix}
& 5 & \\
& -2 & \\
& 3 &
\end{bmatrix} 
& \\
& 
\begin{bmatrix}
& 0 & 0 & 1 &
\end{bmatrix}
\cdot
\begin{bmatrix}
& 5 & \\
& -2 & \\
& 3 &
\end{bmatrix} 
&
\end{bmatrix}
\hspace{0.321cm} \overset{}{=}
\]

\[
\hspace{0.321cm} \overset{}{=} \hspace{0.321cm}
\begin{bmatrix}
& (1)(5) + (0)(-2) + (0)(3) & \\
& (0)(5) + (1)(-2) + (0)(3) & \\
& (0)(5) + (0)(-2) + (1)(3) &
\end{bmatrix}
\hspace{0.321cm} {=} \hspace{0.321cm}
\boxed
{
\begin{bmatrix}
& 5 & \\
& -2 & \\
& 3 &
\end{bmatrix}_{3 \times 1}^E
}
\]


\[
iii)\hspace{0.321cm}\textbf{Έχουμε:} \hspace{0.54321cm}
\begin{bmatrix}
& 2 & 0 & \\
& 1 & 3 &
\end{bmatrix}_{2 \times 2}^F
\cdot \hspace{0.35cm}
\begin{bmatrix}
& 1 & \\
& 1 &
\end{bmatrix}_{2 \times 1}^G
\hspace{0.321cm} \overset{*}{=} \hspace{0.321cm}
\begin{bmatrix}
& 
\begin{bmatrix}
& 2 & 0 &
\end{bmatrix}
\cdot
\begin{bmatrix}
& 1 & \\
& 1 &
\end{bmatrix} 
& \\
& 
\begin{bmatrix}
& 1 & 3 &
\end{bmatrix}
\cdot
\begin{bmatrix}
& 1 & \\
& 1 &
\end{bmatrix} 
&
\end{bmatrix}
\hspace{0.321cm} \overset{}{=}
\]

\[
\hspace{0.321cm} \overset{}{=} \hspace{0.321cm}
\begin{bmatrix}
& (2)(1) + (0)(1) & \\
& (1)(1) + (1)(3) &
\end{bmatrix}
\hspace{0.321cm} {=} \hspace{0.321cm}
\boxed
{
\begin{bmatrix}
& 2 & \\
& 4 &
\end{bmatrix}_{2 \times 1}^H
}
\]


% ---------------------------------------------------------------------------------------


\vspace{2.000cm}
{{\textbf{·  \underline{3\textsuperscript{η} Άσκηση:}}}}
\vspace{0.321cm}


\[
\textbf{Έχουμε:} \hspace{0.54321cm}
\left\{
\begin{array}{l}
x_1 -2x_2 - 3x_3 = -1 \\
3x_1 + x_2 + x_3 = 4 \\
11x_1 - x_2 - 3x_3 = 10
\end{array}
\right.
\hspace{0.321cm} \overset{\textbf{Π}_{3\times4}}{\implies} \hspace{0.321cm}
\begin{bmatrix}
& \boxed{1} & -2 & -3 & | & -1 & \\
& 3 & 1 & 1 & | & 4 & \\
& 11 & -1 & -3 & | & 10 &
\end{bmatrix}
\hspace{0.321cm} \overset{r_2=r_2-\lambda r_1}{\underset{\lambda=\frac{3}{1}=3}{\implies}} \hspace{0.321cm}
\begin{bmatrix}
& \boxed{1} & -2 & -3  & -1 & \\
& \underline{ 0 } & 7 & 10 &  7 & \\
& 11 & -1 & -3 & 10 &
\end{bmatrix} \hspace{0.75cm}
\]

\[
\overset{r_3=r_3-\lambda r_1}{\underset{\lambda=\frac{11}{1}=11}{\implies}} \hspace{0.321cm}
\begin{bmatrix}
& 1 & -2 & -3  & -1 & \\
& \underline{ 0 } & \boxed{7} & 10 & 7 & \\
& \underline{0} & 21 & 30 & 21 &
\end{bmatrix}
\hspace{0.321cm} \overset{r_3=r_3-\lambda r_2}{\underset{\lambda=21/7}{\implies}} \hspace{0.321cm}
\begin{bmatrix}
& 1 & -2 & -3  & -1 & \\
& \underline{ 0 } & 7 & 10 &  7 & \\
& \underline{0} & \underline{0} & \boxed{0} & 0 &
\end{bmatrix}
\hspace{0.321cm} {\implies} \hspace{0.321cm}
\begin{bmatrix}
& 1 & -2 & -3 & | & -1 & \\
& 0 & 7 & 10 & | & 7 & \\
& 0 & 0 & 0 & | & 0 &
\end{bmatrix}
\]

\[
\hspace{0.321cm} {\implies} \hspace{0.321cm}
\left\{
\begin{array}{l}
1x_1 - 2x_2 - 3x_3 = -1 \\
0x_1 + 7x_2 + 10x_3 = 7 \\
0x_1 \pm 0x_2 \pm 0x_3 = 0
\end{array}
\right.
\hspace{0.321cm} \overset{}{\implies} \hspace{0.321cm}
\left\{
\begin{array}{l}
x_1 - 2x_2 - 3x_3 = -1 \\
7x_2 + 10x_3 = 7 \\
0 = 0
\end{array}
\right.
\]

\textbf{(a) Άπειρες λύσεις (αόριστο σύστημα) ...}

\vspace{1.2345cm}

\textbf{(b) Ένα σύστημα ονομάζεται ιδιόμορφο \textlatin{(singular)} όταν ο πίνακας συντελεστών του έχει ορίζουσα που ισούτε με μηδέν. Αυτό σημαίνει ότι οι γραμμές ή οι στήλες του πίνακα είναι γραμμικά εξαρτημένες, και το σύστημα είτε δεν έχει λύση είτε έχει άπειρες λύσεις.}

\textbf{Για παράδειγμα:} \hspace{0.123cm} \textlatin{x+y=6} \hspace{0.123cm} \textbf{και} \hspace{0.123cm} \textlatin{2x+2y=12} .

\vspace{0.7654321cm}
\newpage

% ---------------------------------------------------------------------------------------

\vspace{1.000cm}
{{\textbf{·  \underline{5\textsuperscript{η} Άσκηση:}}}}
\vspace{0.321cm}


\[
(\alpha) \hspace{0.321cm} \textbf{Έχουμε:} \hspace{0.54321cm}
\begin{bmatrix}
& 1 & 1 & 1 & \\
& 1 & 2 & 2 & \\
& 5 & 7 & 8 &
\end{bmatrix}_{3 \times 3}^{A}
\hspace{0.321cm} \overset{*}{\underset{}{\implies}} \hspace{0.321cm}
\begin{bmatrix}
& \boxed{1} & 1 & 1 & | & 1 & 0 & 0 & \\
& 1 & 2 & 2 & | & 0 & 1 & 0 & \\
& 5 & 7 & 8 & | & 0 & 0 & 1 & 
\end{bmatrix}^{Α|Ι}_{3 \times 6}
\hspace{0.321cm} {\implies} \hspace{0.321cm}
\]

\[
\overset{r_2=r_2-\lambda r_1}{\underset{\lambda=1/1=1}{\implies}} \hspace{0.321cm}
\begin{bmatrix}
& \boxed{1} & 1 & 1 & | & 1 & 0 & 0 & \\
& \underline{0} & 1 & 1 & | & -1 & 1 & 0 & \\
& 5 & 7 & 8 & | & 0 & 0 & 1 & 
\end{bmatrix}
\hspace{0.321cm} \overset{r_3=r_3-\lambda r_1}{\underset{\lambda=5/1=5}{\implies}} \hspace{0.321cm}
\begin{bmatrix}
& 1 & 1 & 1 & | & 1 & 0 & 0 & \\
& \underline{0} & \boxed{1} & 1 & | & -1 & 1 & 0 & \\
& \underline{0} & 2 & 3 & | & -5 & 0 & 1 & 
\end{bmatrix}
\hspace{0.321cm} \overset{r_3=r_3-\lambda r_2}{\underset{\lambda=2/1=2}{\implies}} \hspace{0.321cm}
\]

\[
\hspace{0.321cm} \overset{r_3=r_3-\lambda r_2}{\underset{\lambda=2/1=2}{\implies}} \hspace{0.321cm}
\begin{bmatrix}
& 1 & 1 & 1 & | & 1 & 0 & 0 & \\
& \underline{0} & 1 & 1 & | & -1 & 1 & 0 & \\
& \underline{0} & \underline{0} & \boxed{1} & | & -3 & -2 & 1 & 
\end{bmatrix}
\hspace{0.321cm} \implies \hspace{0.321cm}
\begin{bmatrix}
& 1 & 1 & 1 & | & 1 & 0 & 0 & \\
& 0 & 1 & 1 & | & -1 & 1 & 0 & \\
& 0 & 0 & 1 & | & -3 & -2 & 1 & 
\end{bmatrix}
\hspace{0.321cm} \implies \hspace{0.321cm}
\]

\[
\hspace{0.321cm} \overset{r_2=r_2-\lambda r_3}{\underset{\lambda=1}{\implies}} \hspace{0.321cm}
\begin{bmatrix}
& 1 & 1 & 1 & | & 1 & 0 & 0 & \\
& 0 & 1 & \underline{0} & | & 2 & 3 & -1 & \\
& 0 & 0 & 1 & | & -3 & -2 & 1 & 
\end{bmatrix}
\hspace{0.321cm} \overset{r_1=r_1-\lambda r_3}{\underset{\lambda=1}{\implies}} \hspace{0.321cm}
\begin{bmatrix}
& 1 & 1 & \underline{0} & | & 4 & 2 & -1 & \\
& 0 & 1 & \underline{0} & | & 2 & 3 & -1 & \\
& 0 & 0 & 1 & | & -3 & -2 & 1 & 
\end{bmatrix}
\hspace{0.321cm} \overset{r_1=r_1-\lambda r_2}{\underset{\lambda=1}{\implies}} \hspace{0.321cm}
\]

\[
\hspace{0.321cm} \overset{r_1=r_1-\lambda r_2}{\underset{\lambda=1}{\implies}} \hspace{0.321cm}
\begin{bmatrix}
& 1 & \underline{0} & \underline{0} & | & 2 & -1 & 0 & \\
& 0 & 1 & \underline{0} & | & 2 & 3 & -1 & \\
& 0 & 0 & 1 & | & -3 & -2 & 1 & 
\end{bmatrix}^{I|A^{-1}}_{3 \times 6}
\hspace{0.321cm} \overset{... \hspace{0.150cm}}{\implies} \hspace{0.321cm}
\boxed
{
\begin{bmatrix}
& 3 & -1 & 0 & \\
& 1 & 3 & -1 & \\
& -3 & -2 & 1 & 
\end{bmatrix}_{3 \times 3}^{Α^{-1}}
}
\]

\vspace{1cm}

\[
(\beta) \hspace{0.321cm} \textbf{Έχουμε:} \hspace{0.54321cm}
\begin{bmatrix}
& 6 & -2 & 1 & \\
& -4 & 1 & -1 & \\
& 1 & 0 & 1 &
\end{bmatrix}_{3 \times 3}^{B}
\hspace{0.321cm} \overset{*}{\underset{}{\implies}} \hspace{0.321cm}
\begin{bmatrix}
& \boxed{6} & -2 & 1 & | & 1 & 0 & 0 & \\
& -4 & 1 & -1 & | & 0 & 1 & 0 & \\
& 1 & 0 & 1 & | & 0 & 0 & 1 & 
\end{bmatrix}^{B|Ι}_{3 \times 6}
\hspace{0.321cm} {\implies} \hspace{0.321cm}
\]

\[
\overset{r_2=r_2-\lambda r_1}{\underset{\lambda=-4/6=-2/3}{\implies}} \hspace{0.321cm}
\begin{bmatrix}
& \boxed{6} & -2 & 1 & | & 1 & 0 & 0 & \\
& \underline{0} & -\frac{1}{3} & -\frac{1}{3} & | & \frac{2}{3} & 1 & 0 & \\
& 1 & 0 & 1 & | & 0 & 0 & 1 & 
\end{bmatrix}
\hspace{0.321cm} \overset{r_3=r_3-\lambda r_1}{\underset{\lambda=1/6}{\implies}} \hspace{0.321cm}
\begin{bmatrix}
& 6 & -2 & 1 & | & 1 & 0 & 0 & \\
& \underline{0} & \boxed{-\frac{1}{3}} & -\frac{1}{3} & | & \frac{2}{3} & 1 & 0 & \\
& \underline{0} & -\frac{1}{3} & \frac{5}{6} & | & -\frac{1}{6} & 0 & 1 & 
\end{bmatrix}
\hspace{0.321cm} \overset{r_3=r_3-\lambda r_2}{\underset{\lambda=1...}{\implies}} \hspace{0.321cm}
\]

\[
\hspace{0.321cm} \overset{r_3=r_3-\lambda r_2}{\underset{\lambda=1...}{\implies}} \hspace{0.321cm}
\begin{bmatrix}
& 6 & -2 & 1 & | & 1 & 0 & 0 & \\
& \underline{0} & -\frac{1}{3} & -\frac{1}{3} & | & \frac{2}{3} & 1 & 0 & \\
& \underline{0} & \underline{0} & \boxed{\frac{7}{6}} & | & -\frac{5}{6} & -1 & 1 & 
\end{bmatrix}
\hspace{0.321cm} \implies \hspace{0.321cm}
\begin{bmatrix}
& 6 & -2 & 1 & | & 1 & 0 & 0 & \\
& 0 & -1/3 & -1/3 & | & 2/3 & 1 & 0 & \\
& 0 & 0 & 7/6 & | & -5/6 & -1 & 1 & 
\end{bmatrix}
\hspace{0.321cm} \implies \hspace{0.321cm}
\]

\[
\hspace{0.321cm} \overset{r_3 = r_3 \cdot \frac{6}{7}}{\underset{}{\implies}} \hspace{0.321cm}
\begin{bmatrix}
& 6 & -2 & 1 & | & 1 & 0 & 0 & \\
& 0 & -\frac{1}{3} & -\frac{1}{3} & | & \frac{2}{3} & 1 & 0 & \\
& 0 & 0 & \underline{1} & | & -\frac{5}{7} & -\frac{6}{7} & \frac{6}{7} & 
\end{bmatrix}
\hspace{0.321cm} \overset{r_2=r_2 \cdot (-3)}{\underset{}{\implies}} \hspace{0.321cm}
\begin{bmatrix}
& 6 & -2 & 1 & | & 1 & 0 & 0 & \\
& 0 & \underline{1} & \underline{1} & | & -2 & -3 & 0 & \\
& 0 & 0 & \underline{1} & | & -\frac{5}{7} & -\frac{6}{7} & \frac{6}{7} & 
\end{bmatrix}
\hspace{0.321cm} \overset{r_2=r_2-\lambda r_3}{\underset{\lambda=1}{\implies}} \hspace{0.321cm}
\]

\[
\hspace{0.321cm} \overset{r_2=r_2-\lambda r_3}{\underset{\lambda=1}{\implies}} \hspace{0.321cm}
\begin{bmatrix}
& 6 & -2 & 1 & | & 1 & 0 & 0 & \\
& 0 & \underline{1} & \underline{0} & | & -\frac{9}{7} & -\frac{15}{7} & -\frac{6}{7} & \\
& 0 & 0 & \underline{1} & | & -\frac{5}{7} & -\frac{6}{7} & \frac{6}{7} & 
\end{bmatrix}
\hspace{0.321cm} \overset{r_1=r_1 \cdot \frac{1}{6}}{\underset{}{\implies}} \hspace{0.321cm}
\begin{bmatrix}
& \underline{1} & -\frac{1}{3} & \frac{1}{6} & | & \frac{1}{6} & 0 & 0 & \\
& 0 & \underline{1} & \underline{0} & | & -\frac{9}{7} & -\frac{15}{7} & -\frac{6}{7} & \\
& 0 & 0 & \underline{1} & | & -\frac{5}{7} & -\frac{6}{7} & \frac{6}{7} & 
\end{bmatrix}
\hspace{0.321cm} \overset{r_1=r_1-\lambda r_3}{\underset{\lambda=1/6}{\implies}} \hspace{0.321cm}
\]

\[
\hspace{0.321cm} \overset{r_1=r_1-\lambda r_3}{\underset{\lambda=1/6}{\implies}} \hspace{0.321cm}
\begin{bmatrix}
& \underline{1} & -\frac{1}{3} & \underline{0} & | & \frac{2}{7} & \frac{1}{7} & -\frac{1}{7} & \\
& 0 & \underline{1} & \underline{0} & | & -\frac{9}{7} & -\frac{15}{7} & -\frac{6}{7} & \\
& 0 & 0 & \underline{1} & | & -\frac{5}{7} & -\frac{6}{7} & \frac{6}{7} & 
\end{bmatrix}
\hspace{0.321cm} \overset{r_1=r_1-\lambda r_2}{\underset{\lambda=-1/3}{\implies}} \hspace{0.321cm}
\begin{bmatrix}
& \underline{1} & \underline{0} & \underline{0} & | & -\frac{1}{7} & -\frac{4}{7} & -\frac{6}{7} & \\
& 0 & \underline{1} & \underline{0} & | & -\frac{9}{7} & -\frac{15}{7} & -\frac{6}{7} & \\
& 0 & 0 & \underline{1} & | & -\frac{5}{7} & -\frac{6}{7} & \frac{6}{7} & 
\end{bmatrix}_{3 \times 6}^{I \hspace{0.1cm} | \hspace{0.1cm} B^{-1}}
\hspace{0.321cm} \overset{}{\underset{}{\implies}} \hspace{0.321cm}
\]

\[
\hspace{0.321cm} \overset{... \hspace{0.150cm}}{\implies} \hspace{0.321cm}
\boxed
{
\begin{bmatrix}
& -\frac{1}{7} & -\frac{4}{7} & -\frac{6}{7} & \\
& -\frac{9}{7} & -\frac{15}{7} & -\frac{6}{7} & \\
& -\frac{5}{7} & -\frac{6}{7} & \frac{6}{7} & 
\end{bmatrix}_{3 \times 3}^{B^{-1}}
}
\hspace{0.321cm} \overset{... \hspace{0.150cm}}{\implies} \hspace{0.321cm}
{
\begin{bmatrix}
& - {1}/{7} & - {4}/{7} & - {6}/{7} & \\
& - {9}/{7} & - {15}/{7} & - {6}/{7} & \\
& - {5}/{7} & - {6}/{7} & + {6}/{7} & 
\end{bmatrix}_{3 \times 3}^{B^{-1}}
}
\]


% ---------------------------------------------------------------------------------------


\newpage

\vspace{1.000cm}
{{\textbf{·  \underline{6\textsuperscript{η} Άσκηση:}}}}
\vspace{0.321cm}

\[
\textbf{Α.} \hspace{0.321cm} \textbf{Εάν το μοναδικό στοιχείο του είναι διάφορο του μηδενός.}
\]

\[
\textbf{Β.} \hspace{0.321cm} \textbf{Μόνο όταν η ορίζουσα του είναι διάφορη του μηδενός.}
\]

\[
\textbf{Γ.} \hspace{0.321cm} \textbf{Μόνο όταν όλα τα στοιχεία της διαγωνίου του είναι διάφορα του μηδενός.}
\]


% ---------------------------------------------------------------------------------------


\vspace{1.2345cm}
{{\textbf{·  \underline{7\textsuperscript{η} Άσκηση:}}}}
\vspace{0.321cm}

\[
(\alpha) \hspace{0.1cm}  \textbf{Είναι: } \hspace{0.321cm} 
AB = AC \hspace{0.1cm} \Rightarrow \hspace{0.1cm}
A^{-1}(AB) = A^{-1}(AC) \hspace{0.1cm} \Rightarrow \hspace{0.1cm} 
(A^{-1}A)B = (A^{-1}A)C \hspace{0.1cm} \Rightarrow \hspace{0.1cm}
IB = IC \hspace{0.1cm} \Rightarrow \hspace{0.1cm} \boxed{B=C}
\]

\[
(\beta) \hspace{0.1cm}  \textbf{Ένα \hspace{0.1cm}π.χ. είναι: } \hspace{0.321cm}
\boxed
{
B_{2 \times 2} \hspace{0.1cm} = \hspace{0.1cm}
\begin{bmatrix}
& 0 & 1 & \\
& 2 & 3 &
\end{bmatrix}
}
\hspace{0.321cm} \underline{\textbf{και}} \hspace{0.321cm}
\boxed
{
C_{2 \times 2} \hspace{0.1cm} = \hspace{0.1cm}
\begin{bmatrix}
& 0 & 1 & \\
& 3 & 2 &
\end{bmatrix}
}
\]


% ---------------------------------------------------------------------------------------


\vspace{1.2345cm}
{{\textbf{·  \underline{8\textsuperscript{η} Άσκηση:}}}}
\vspace{0.321cm}

\[
\textbf{Έχουμε:} \hspace{0.321cm}
A_{3 \times 2} \hspace{0.1cm} = \hspace{0.1cm}
\begin{bmatrix}
& 2 & 1 & \\
& 1 & 2 & \\
& 3 & 4 &
\end{bmatrix}
\hspace{0.321cm} \textbf{και} \hspace{0.321cm}
B_{3 \times 2} \hspace{0.1cm} = \hspace{0.1cm}
\begin{bmatrix}
& 1 & 3 & \\
& 4 & 2 & \\
& 5 & 6 &
\end{bmatrix}
\]

\[
\textbf{Αρχικά:} \hspace{0.321cm}
A^{\textbf{τ}}_{2 \times 3} \hspace{0.1cm} = \hspace{0.1cm}
\begin{bmatrix}
& 2 & 1 & 3 & \\
& 1 & 2 & 4 & 
\end{bmatrix}
\hspace{0.321cm} \textbf{και} \hspace{0.321cm}
B^{\textbf{τ}}_{2 \times 3} \hspace{0.1cm} = \hspace{0.1cm}
\begin{bmatrix}
& 1 & 4 & 5 & \\
& 3 & 2 & 6 & 
\end{bmatrix}
\hspace{0.321cm}
\textbf{(αλλαγή γραμμών με στήλλες ...)}
\]

\[
\alpha) \hspace{0.321cm} A^{\textbf{τ}}_{2 \times 3} \cdot B_{3 \times 2} = C^1_{2 \times 2} = 
\begin{bmatrix}
& (2)(1)+(1)(4)+(3)(5) & (2)(3)+(1)(2)+(3)(6) & \\
& (1)(1)+(2)(4)+(4)(5) & (1)(3)+(2)(2)+(4)(6) & 
\end{bmatrix}
=
\boxed
{
\begin{bmatrix}
& 21 & 26 & \\
& 29 & 31 & 
\end{bmatrix}
}
\]

\[
\beta) \hspace{0.321cm} B^{\textbf{τ}}_{2 \times 3} \cdot A_{3 \times 2} = C^2_{2 \times 2} = 
\begin{bmatrix}
& (1)(2)+(4)(1)+(5)(3) & (1)(1)+(4)(2)+(5)(4) & \\
& (3)(2)+(2)(1)+(6)(3) & (3)(1)+(2)(2)+(6)(4) & 
\end{bmatrix}
=
\boxed
{
\begin{bmatrix}
& 21 & 29 & \\
& 26 & 31 & 
\end{bmatrix}
}
\]

\[
\gamma) \hspace{0.321cm} A_{3 \times 2} \cdot B^{\textbf{τ}}_{2 \times 3} = C^3_{3 \times 3} = 
\begin{bmatrix}
& (2)(1)+(1)(3) & (2)(4)+(1)(2) & (2)(5)+(1)(6) & \\
& (1)(1)+(2)(3) & (1)(4)+(2)(2) & (1)(5)+(2)(6) & \\
& (3)(1)+(4)(3) & (3)(4)+(4)(2) & (3)(5)+(4)(6) &
\end{bmatrix}
=
\boxed 
{
\begin{bmatrix}
& 5 & 10 & 16 & \\
& 7 & 8 & 17 & \\
& 15 & 20 & 39 &
\end{bmatrix}
}
\hspace{0.11111cm}
\]

\[
\delta) \hspace{0.321cm} B_{3 \times 2} \cdot A^{\textbf{τ}}_{2 \times 3} = C^4_{3 \times 3} = 
\begin{bmatrix}
& (1)(2)+(3)(1) & (1)(1)+(3)(2) & (1)(3)+(3)(4) & \\
& (4)(2)+(2)(1) & (4)(1)+(2)(2) & (4)(3)+(2)(4) & \\
& (5)(2)+(6)(1) & (5)(1)+(6)(2) & (5)(3)+(6)(4) &
\end{bmatrix}
=
\boxed 
{
\begin{bmatrix}
& 5 & 7 & 15 & \\
& 10 & 8 & 20 & \\
& 16 & 17 & 39 &
\end{bmatrix}
}
\hspace{0.11111cm}
\]


% ---------------------------------------------------------------------------------------


\newpage
\vspace{1.2345cm}
{{\textbf{·  \underline{10\textsuperscript{η} Άσκηση:}}}}
\vspace{1   cm}


\[
\alpha) \hspace{0.321cm} \textbf{Έχουμε:} \hspace{0.54321cm}
A^1_{2 \times 2} \hspace{0.321cm} = \hspace{0.321cm}
\begin{bmatrix}
& \boxed{3}  & 2 & \\
& 12 & 9 &
\end{bmatrix}
\hspace{0.321cm} \overset{r_2=r_2-\lambda r_1}{\underset{\underline{ \lambda=4^* }}{\implies}} \hspace{0.321cm}
\begin{bmatrix}
& 3  & 2 & \\
& 0 & \boxed{1} &
\end{bmatrix}
\hspace{0.321cm} \implies \hspace{0.321cm}
\begin{bmatrix}
& 3  & 2 & \\
& 0 & 1 &
\end{bmatrix}^*
\]

\[
\textbf{Άρα:} \hspace{0.321cm}
\boxed
{
U_{2 \times 2} = 
\begin{bmatrix}
& 3  & 2 & \\
& 0 & 1 &
\end{bmatrix}
\hspace{0.321cm} \textbf{και} \hspace{0.321cm} 
L_{2 \times 2} = 
\begin{bmatrix}
& 1 & 0 & \\
& 4 & 1 &
\end{bmatrix}
}
\]

\vspace{1.1cm}

\[
\beta) \hspace{0.321cm} \textbf{Έχουμε:} \hspace{0.54321cm}
A^2_{3 \times 3} \hspace{0.321cm} = \hspace{0.321cm}
\begin{bmatrix}
& \boxed{2} & 1 & 1 & \\
& 1 & 2 & 1 & \\
& 1 & 1 & 2 &
\end{bmatrix}
\hspace{0.321cm} 
\overset{r_2=r_2-\lambda r_1}{\underset{\underline{ \lambda=1/2 }}{\implies}} 
\hspace{0.321cm}
\begin{bmatrix}
& \boxed{2} & 1 & 1 & \\
& 0 & \frac{3}{2} & \frac{1}{2} & \\
& 1 & 1 & 2 &
\end{bmatrix}
\hspace{0.321cm} 
\overset{r_3=r_3-\lambda r_1}{\underset{\underline{ \lambda=1/2 }}{\implies}} 
\hspace{0.321cm}
\begin{bmatrix}
& 2 & 1 & 1 & \\
& 0 & \boxed{\frac{3}{2}} & \frac{1}{2} & \\
& 0 & \frac{1}{2} & \frac{3}{2} &
\end{bmatrix}
\]

\[
\overset{r_3=r_3-\lambda r_2}{\underset{\underline{ \lambda=1/3 }}{\implies}} 
\begin{bmatrix}
& 2 & 1 & 1 & \\
& 0 & {\frac{3}{2}} & \frac{1}{2} & \\
& 0 & 0 & \boxed{\frac{4}{3}} &
\end{bmatrix}
\hspace{0.321cm} \implies \hspace{0.321cm}
\begin{bmatrix}
& 2 & 1 & 1 & \\
& 0 & 3/2 & 1/2 & \\
& 0 & 0 & 4/3 &
\end{bmatrix}
\]

\[
\textbf{Άρα:} \hspace{0.321cm}
\boxed
{
U_{3 \times 3} = 
\begin{bmatrix}
& 2 & 1 & 1 & \\
& 0 & 3/2 & 1/2 & \\
& 0 & 0 & 4/3 &
\end{bmatrix}
\hspace{0.321cm} \textbf{και} \hspace{0.321cm} 
L_{3 \times 3} = 
\begin{bmatrix}
& 1 & 0 & 0 & \\
& 1/2 & 1 & 0 & \\
& 1/2 & 1/3 & 1 &
\end{bmatrix}
}
\]

\vspace{1.1cm}

\[
\gamma) \hspace{0.321cm} \textbf{Έχουμε:} \hspace{0.54321cm}
A^3_{3 \times 3} \hspace{0.321cm} = \hspace{0.321cm}
\begin{bmatrix}
& \boxed{1} & 1 & 1 & \\
& 1 & 3 & 3 & \\
& 1 & 3 & 9 &
\end{bmatrix}
\hspace{0.321cm} 
\overset{r_2=r_2-\lambda r_1}{\underset{\underline{ \lambda=1 }}{\implies}} 
\hspace{0.321cm}
\begin{bmatrix}
& \boxed{1} & 1 & 1 & \\
& 0 & 2 & 2 & \\
& 1 & 3 & 9 &
\end{bmatrix}
\hspace{0.321cm} 
\overset{r_3=r_3-\lambda r_1}{\underset{\underline{ \lambda=1 }}{\implies}} 
\hspace{0.321cm}
\begin{bmatrix}
& 1 & 1 & 1 & \\
& 0 & \boxed{2} & 2 & \\
& 0 & 2 & 8 &
\end{bmatrix}
\]

\[
\overset{r_3=r_3-\lambda r_2}{\underset{\underline{ \lambda=1 }}{\implies}} 
\begin{bmatrix}
& 1 & 1 & 1 & \\
& 0 & 2 & 2 & \\
& 0 & 0 & \boxed{6} &
\end{bmatrix}
\hspace{0.321cm} \implies \hspace{0.321cm}
\begin{bmatrix}
& 1 & 1 & 1 & \\
& 0 & 2 & 2 & \\
& 0 & 0 & 6 &
\end{bmatrix}
\]

\[
\textbf{Άρα:} \hspace{0.321cm}
\boxed
{
U_{3 \times 3} = 
\begin{bmatrix}
& 1 & 1 & 1 & \\
& 0 & 2 & 2 & \\
& 0 & 0 & 6 &
\end{bmatrix}
\hspace{0.321cm} \textbf{και} \hspace{0.321cm} 
L_{3 \times 3} = 
\begin{bmatrix}
& 1 & 0 & 0 & \\
& 1 & 1 & 0 & \\
& 1 & 1 & 1 &
\end{bmatrix}
}
\]


% ---------------------------------------------------------------------------------------


\newpage
\vspace{1.2345cm}
{{\textbf{·  \underline{11\textsuperscript{η} Άσκηση:}}}}
\vspace{1cm}


\[
\alpha) \hspace{0.321cm} \textbf{Έχουμε:} \hspace{0.54321cm}
A^{\alpha}_{3 \times 3} \hspace{0.321cm} = \hspace{0.321cm}
\begin{bmatrix}
& \boxed{0}^* & 1 & 1 & \\
& 1 & 0 & 1 & \\
& 4 & 5 & 6 &
\end{bmatrix}
\hspace{0.321cm} 
\overset{r_1 \leftarrow r_3}{\underset{{r_3 \leftarrow r_1}}{\implies}} 
\hspace{0.321cm}
\begin{bmatrix}
& \boxed{4} & 5 & 6 & \\
& 1 & 0 & 1 & \\
& 0 & 1 & 1 &
\end{bmatrix}
\hspace{0.321cm} \textbf{και έτσι} \hspace{0.321cm}
P_{3 \times 3} \hspace{0.321cm} = \hspace{0.321cm}
\begin{bmatrix}
& 0 & 0 & 1 & \\
& 0 & 1 & 0 & \\
& 1 & 0 & 0 &
\end{bmatrix}
\hspace{0.25cm}
\]

\[
\overset{r_2 = r_2 - \lambda r_1}{\underset{\underline{\lambda = 1/4}}{\implies}} 
\hspace{0.321cm}
\begin{bmatrix}
& 4 & 5 & 6 & \\
& 0 & \boxed{-\frac{5}{4}} & -\frac{1}{2} & \\
& 0 & 1 & 1 &
\end{bmatrix}
\hspace{0.321cm}
\overset{r_3 = r_3 - \lambda r_2}{\underset{\underline{\lambda = - 4/5}}{\implies}} 
\hspace{0.321cm}
\begin{bmatrix}
& 4 & 5 & 6 & \\
& 0 & -\frac{5}{4} & -\frac{1}{2} & \\
& 0 & 0 & \boxed{\frac{3}{5}} &
\end{bmatrix}
\hspace{0.321cm}
\overset{... \hspace{0.1cm}}{\implies}
\hspace{0.321cm}
\begin{bmatrix}
& 4 & 5 & 6 & \\
& 0 & -5/4 & -1/2 & \\
& 0 & 0 & 3/5 &
\end{bmatrix}
\]

\[
\textbf{Άρα} \hspace{0.321cm}
L_{3 \times 3} \hspace{0.321cm} = \hspace{0.321cm}
\begin{bmatrix}
& 1 & 0 & 0 & \\
& 1/4 & 1 & 0 & \\
& 0 & -4/5 & 1 &
\end{bmatrix}
\hspace{0.321cm} \textbf{και} \hspace{0.321cm}
D_{3 \times 3} \hspace{0.321cm} = \hspace{0.321cm}
\begin{bmatrix}
& 4 & 0 & 0 & \\
& 0 & -5/4 & 0 & \\
& 0 & 0 & 3/5 &
\end{bmatrix}
\hspace{0.321cm} \textbf{και} \hspace{0.321cm}
U_{3 \times 3} \hspace{0.321cm} = \hspace{0.321cm}
\begin{bmatrix}
& 1 & 5/4 & 3/2 & \\
& 0 & 1 & 2/5 & \\
& 0 & 0 & 1 &
\end{bmatrix}
\hspace{0.25cm}
\]

\[
\textbf{Τελικά:} \hspace{0.321cm}
\boxed
{
\begin{bmatrix}
& 0 & 0 & 1 & \\
& 0 & 1 & 0 & \\
& 1 & 0 & 0 &
\end{bmatrix}^P
\cdot
\begin{bmatrix}
& 0 & 1 & 1 & \\
& 1 & 0 & 1 & \\
& 4 & 5 & 6 &
\end{bmatrix}^A
=
\begin{bmatrix}
& 1 & 0 & 0 & \\
& 1/4 & 1 & 0 & \\
& 0 & -4/5 & 1 &
\end{bmatrix}^L
\cdot
\begin{bmatrix}
& 4 & 0 & 0 & \\
& 0 & -5/4 & 0 & \\
& 0 & 0 & 3/5 &
\end{bmatrix}^D
\cdot
\begin{bmatrix}
& 1 & 5/4 & 3/2 & \\
& 0 & 1 & 2/5 & \\
& 0 & 0 & 1 &
\end{bmatrix}^U
}
\hspace{0.54321cm}
\]



\vspace{2cm}



\[
\beta) \hspace{0.321cm} \textbf{Έχουμε:} \hspace{0.54321cm}
A^{\beta}_{3 \times 3} \hspace{0.321cm} = \hspace{0.321cm}
\begin{bmatrix}
& \boxed{3} & 6 & 3 & \\
& 6 & 12 & 6 & \\
& 3 & 3 & 3 &
\end{bmatrix}
\hspace{0.321cm} 
\overset{r_2 = r_2 - \lambda r_1}{\underset{\underline{\lambda_{21} = 2}}{\implies}}
\hspace{0.321cm} 
\begin{bmatrix}
& \boxed{3} & 6 & 3 & \\
& 0 & 0 & 0 & \\
& 3 & 3 & 3 &
\end{bmatrix}
\hspace{0.321cm} 
\overset{r_3 = r_3 - \lambda r_1}{\underset{\underline{\lambda_{31} = 1}}{\implies}}
\hspace{0.321cm} 
\begin{bmatrix}
& 3 & 6 & 3 & \\
& 0 & \boxed{0}^* & 0 & \\
& 0 & -3 & 0 &
\end{bmatrix}
\]

\[
\hspace{0.321cm} 
\overset{r_2 \leftarrow r_3}{\underset{{r_3 \leftarrow r_2}}{\implies}} 
\hspace{0.321cm}
\begin{bmatrix}
& 3 & 6 & 3 & \\
& 0 & \boxed{-3} & 0 & \\
& 0 & 0 & 0 &
\end{bmatrix}
\hspace{0.321cm} \implies \hspace{0.321cm}
\begin{bmatrix}
& 3 & 6 & 3 & \\
& 0 & -3 & 0 & \\
& 0 & 0 & \boxed{0}^! &
\end{bmatrix}
\hspace{0.321cm} \overset{^?}{\underset{...}{\implies}} \hspace{0.321cm}
\begin{bmatrix}
& 3 & 6 & 3 & \\
& 0 & -3 & 0 & \\
& 0 & 0 & 0 &
\end{bmatrix}^{***}
\]


\[
\textbf{Αφού ο οδηγός είνια μηδέν και δεν γίνονται αντιμεταθέσεις, η παραγωντοποίηση δεν ολοκληρώνεται !} \hspace{0.001cm}
\]

\[
\textbf{Τελικά:} \hspace{0.321cm}
\boxed
{
\begin{bmatrix}
& 1 & 0 & 0 & \\
& 0 & 0 & 1 & \\
& 0 & 1 & 0 &
\end{bmatrix}^P
\cdot
\begin{bmatrix}
& 3 & 6 & 3 & \\
& 6 & 12 & 6 & \\
& 3 & 3 & 3 &
\end{bmatrix}^A
=
\begin{bmatrix}
& 1 & 0 & 0 & \\
& 2 & 1 & 0 & \\
& 1 & 0 & 1 &
\end{bmatrix}^L
\cdot
\begin{bmatrix}
&  &  &  & \\
&  & \boxed{?} &  & \\
&  &  &  &
\end{bmatrix}^D
\cdot
\begin{bmatrix}
& 3 & 6 & 3 & \\
& 0 & -3 & 0 & \\
& 0 & 0 & 0 &
\end{bmatrix}^U
}
\hspace{0.54321cm}
\]

\[
\textbf{* Υπόθεση:} \hspace{0.321cm}
Δ_{3 \times 3} \hspace{0.25cm} = \hspace{0.25cm}
\begin{bmatrix}
& 3 & 0 & 0 & \\
& 0 & -3 & 0 & \\
& 0 & 0 & 0 \hspace{0.05cm} ? &
\end{bmatrix}
\]

% ---------------------------------------------------------------------------------------

\vspace{3.5cm}
\begin{center}
\boxed{\underline{\textbf{Ασκήσεις 4 και 9 είναι προγραμματιστικές ..!}}}
\end{center}

% ---------------------------------------------------------------------------------------

\newpage

\begin{center}
    \vspace*{\fill}
    \setlength{\fboxsep}{10pt}
    \boxed{\textbf{\LARGE{ΤΕΛΟΣ}}}
    \vspace*{\fill}
\end{center}

\end{document}

% ---------------------------------------------------------------------------------------

% The end ...